\documentclass[
  aps,
  prd,
  twocolumn,
  nofootinbib,
  superscriptaddress,
  longbibliography
]{revtex4-2}

\usepackage{amsmath,amssymb,amsfonts}
\usepackage{bm}
\usepackage{graphicx}
\usepackage{hyperref}

\begin{document}

\title{Coordinate-Time and Proper-Time Brachistochrones\\
in Tolman--Oppenheimer--Volkoff Stellar Interiors}

\author{Iman Rosignoli}
\email{iman.rosignoli01@universitadipavia.it}
\affiliation{Department of Physics, University of Pavia, Pavia, Italy}

\date{\today}

\begin{abstract}
We formulate the relativistic brachistochrone problem inside static,
spherically symmetric stellar interiors generated by the
Tolman--Oppenheimer--Volkoff equations for a prescribed equation of
state. The stellar geometry is determined by the pressure,
energy-density, mass and lapse profiles obtained from the TOV system.
For a massive test particle released from rest at the stellar surface,
we derive two distinct variational problems: the minimization of the
coordinate travel time $t$ measured by an observer at infinity, and the
minimization of the proper time $\tau$ measured along the particle
worldline. Both problems reduce to Fermat-type variational principles
on the spatial geometry of the static slice, with two different
effective refractive indices. The cyclic angular coordinate yields generalized Snell laws and
integral formulae for the opening angle, the turning radius and the
travel times. A universal weak-field expansion in the compactness
parameter $\mu=r_s/R$ shows that the difference between the
coordinate-time and proper-time turning radii,
$q^{(t)}-q^{(\tau)}$, is determined at $O(\mu)$ for \emph{any}
equation of state by two one-dimensional quadratures and one
derivative, without solving the full TOV system. In the
constant-density limit this reduces to the closed form
$\mu\,q_0(1-q_0^2)^2/4$. These results provide a framework for
comparing relativistic brachistochrone curves in constant-density
interiors, polytropic stars and realistic neutron-star equations of
state.
\end{abstract}

\maketitle

\section{Introduction}

The classical brachistochrone problem asks for the curve along which a
particle moving under gravity travels between two fixed points in the
least time. In Newtonian gravity with a uniform gravitational field,
the answer is a cycloid. In a spherically symmetric gravitational
field, the corresponding ``gravitational train'' problem concerns the
curve through a gravitating body that minimizes the travel time between
two surface points. The Newtonian gravity train and its
variants---including non-uniform and polytropic density
profiles---have been studied extensively~\cite{Klotz2015,GjerlovPesnell2019},
and relativistic generalizations have been considered for
uniform-density spheres up to the black-hole limit~\cite{Parker2017,Seel2018},
where the distinction between coordinate and proper travel time already
becomes essential.

In general relativity the problem becomes richer. The travel time
depends not only on the local velocity gained by the falling particle,
but also on the gravitational redshift and on the spatial curvature of
the interior geometry. For a static spacetime one may formulate the
problem as a Fermat-type variational principle: the coordinate travel
time equals the optical length in a positive-definite Riemannian
(Jacobi, or optical) metric on the spatial slices, as established in
full generality by Perlick~\cite{Perlick1990} for stationary
spacetimes; for massive particles the analogous structure is the
energy-dependent Jacobi metric on the static slices~\cite{Gibbons2016}.
The brachistochrone is therefore a geodesic of this effective metric,
with an effective refractive index determined by the lapse function. A
closely related refraction (Snell-law) formulation of relativistic free
fall, including the Schwarzschild case, was given by
Schmidt~\cite{Schmidt2018}.

Previous relativistic treatments often use analytically tractable
interior metrics, such as the constant-density interior Schwarzschild
solution. In particular, the present work extends the coordinate-time
brachistochrone problem studied for the interior Schwarzschild metric
in Ref.~\cite{Rosignoli2026}, where perturbative NLO and NNLO
corrections to the Newtonian hypocycloid were derived in the
compactness parameter $\mu=r_s/R$ and validated against direct
numerical integration. However, a realistic compact star is not of
constant density. For a neutron star, the internal structure is
determined by an equation of state together with the
Tolman--Oppenheimer--Volkoff equations. This suggests a natural
extension of the relativistic brachistochrone problem: given an
equation of state, solve the TOV system, construct the interior metric,
and determine the coordinate-time and proper-time brachistochrone
curves.

In this work we derive the general equations for this problem. We
consider a static, spherically symmetric TOV stellar interior and a
massive test particle released from rest at the stellar surface. We
derive the coordinate-time brachistochrone, which minimizes the time
$t$ measured by a static observer at infinity, and the proper-time
brachistochrone, which minimizes the worldline time $\tau$ of the
particle itself. The two problems are governed by different effective
indices and therefore, in general, by different minimizing curves. The
central analytic result of the paper (Sec.~\ref{sec:weak_field}) is a
model-independent formula for the leading relativistic separation
between the two brachistochrones: the difference of turning radii
$q^{(t)}-q^{(\tau)}$ is determined at $O(\mu)$ for any equation of
state by two one-dimensional quadratures over the Newtonian density
profile alone, without solving the full TOV system. In the
constant-density limit one recovers the closed-form expression
$\mu\,q_0(1-q_0^2)^2/4$, which serves as a direct analytical
benchmark for numerical TOV integrations.

\section{TOV stellar interior}

We consider a static, spherically symmetric perfect-fluid star. The
stress-energy tensor is
\begin{equation}
T^{\mu\nu}
=
\left(\epsilon+p\right)u^\mu u^\nu
+
p\,g^{\mu\nu},
\end{equation}
where $\epsilon(r)$ is the energy density and $p(r)$ is the pressure.
The line element is written as
\begin{equation}
ds^2
=
-A(r)c^2dt^2
+
B(r)dr^2
+
r^2d\Omega^2,
\label{eq:metric_AB}
\end{equation}
where
\begin{equation}
A(r)=e^{2\Phi(r)},
\qquad
B(r)=
\left(
1-\frac{2Gm(r)}{rc^2}
\right)^{-1}.
\end{equation}

The stellar structure is determined by the TOV
equations~\cite{Tolman1939,OppenheimerVolkoff1939}
\begin{equation}
\frac{dm}{dr}
=
\frac{4\pi r^2}{c^2}\,\epsilon(r),
\label{eq:tov_mass}
\end{equation}
and
\begin{multline}
\frac{dp}{dr}
=
-
\frac{G\epsilon(r)m(r)}{c^2r^2}
\left(
1+\frac{p(r)}{\epsilon(r)}
\right)
\left(
1+\frac{4\pi r^3p(r)}{m(r)c^2}
\right)
\\
\times
\left(
1-\frac{2Gm(r)}{rc^2}
\right)^{-1}.
\label{eq:tov_pressure}
\end{multline}
The lapse function is obtained from
\begin{equation}
\frac{d\Phi}{dr}
=
\frac{
G\left[
m(r)+4\pi r^3p(r)/c^2
\right]
}{
r^2c^2
\left(
1-\dfrac{2Gm(r)}{rc^2}
\right)
}.
\label{eq:tov_lapse}
\end{equation}

An equation of state closes the system. We assume that the EOS can be
written as
\begin{equation}
\epsilon=\epsilon(p),
\end{equation}
either analytically, as in a polytropic model, or numerically, as in a
tabulated neutron-star equation of state. Given a central pressure
$p_c$, the integration starts from
\begin{equation}
m(0)=0,
\qquad
p(0)=p_c,
\end{equation}
and proceeds outward until
\begin{equation}
p(R)=0.
\end{equation}
The surface radius is $R$, and the total gravitational mass is
\begin{equation}
M=m(R).
\end{equation}
At the surface, the interior metric is matched to the exterior
Schwarzschild geometry, giving
\begin{equation}
A(R)
=
1-\frac{2GM}{Rc^2}.
\label{eq:surface_lapse}
\end{equation}
In the following we denote
\begin{equation}
A_0=A(R).
\label{eq:A0_def}
\end{equation}
The endpoints of the brachistochrone are taken to lie on the stellar
surface and to be separated by an angular aperture $\Delta$.

\section{Particle released from rest at the surface}

We restrict the motion to the equatorial plane,
$\theta=\pi/2$,
so that the metric becomes
\begin{equation}
ds^2
=
-A(r)c^2dt^2
+
B(r)dr^2
+
r^2d\phi^2.
\label{eq:equatorial_metric}
\end{equation}
The spatial line element on a static slice is
\begin{equation}
d\ell^2
=
B(r)dr^2+r^2d\phi^2.
\label{eq:spatial_line_element}
\end{equation}

We assume the tunnel is frictionless and exerts only forces normal to
the motion, so that the constraint does no work and the particle's
energy is conserved; the local speed then depends on $r$ alone. For a
massive particle moving in the static spacetime, the conserved specific
energy is
\begin{equation}
\mathcal{E}
=
A(r)c^2\frac{dt}{d\tau}.
\end{equation}
If the particle is released from rest at the surface $r=R$, then
\begin{equation}
\mathcal{E}
=
c^2\sqrt{A_0}.
\label{eq:energy_surface_release}
\end{equation}
The local Lorentz factor measured by static observers is
\begin{equation}
\gamma(r)
=
\frac{\mathcal{E}}{c^2\sqrt{A(r)}}
=
\sqrt{\frac{A_0}{A(r)}}.
\end{equation}
Therefore the local speed satisfies
\begin{equation}
\frac{v^2(r)}{c^2}
=
1-\frac{1}{\gamma^2(r)}
=
1-\frac{A(r)}{A_0}.
\label{eq:local_speed}
\end{equation}

The coordinate-time and proper-time brachistochrone problems differ
because
\begin{equation}
dt
=
\frac{d\ell}{\sqrt{A(r)}\,v(r)},
\label{eq:dt_dl}
\end{equation}
whereas
\begin{equation}
d\tau
=
\frac{d\ell}{\gamma(r)v(r)}.
\label{eq:dtau_dl}
\end{equation}

\section{Coordinate-time brachistochrone}

The coordinate travel time is
\begin{equation}
t[\gamma]
=
\int_\gamma dt
=
\frac{\sqrt{A_0}}{c}
\int_\gamma
\frac{d\ell}{\sqrt{A(r)\left[A_0-A(r)\right]}}.
\end{equation}
The overall factor $\sqrt{A_0}$ is constant and does not affect the
minimizing curve. We therefore define the coordinate-time effective
index
\begin{equation}
n_t(r)
=
\frac{1}{\sqrt{A(r)\left[A_0-A(r)\right]}}.
\label{eq:nt_def}
\end{equation}
The coordinate-time brachistochrone is determined by the variational
principle
\begin{equation}
\delta
\int_\gamma
n_t(r)\,d\ell
=
0.
\label{eq:t_variational}
\end{equation}
In an arbitrary parameter $\lambda$, the functional is
\begin{equation}
t[\gamma]
=
\frac{\sqrt{A_0}}{c}
\int
n_t(r)
\sqrt{
B(r)\dot{r}^2+r^2\dot{\phi}^2
}
\,d\lambda,
\end{equation}
where a dot denotes $d/d\lambda$. The corresponding Lagrangian is
\begin{equation}
\mathcal{L}_t
=
n_t(r)
\sqrt{
B(r)\dot{r}^2+r^2\dot{\phi}^2
}.
\end{equation}
Since $\phi$ is cyclic,
\begin{equation}
\frac{\partial\mathcal{L}_t}{\partial\dot{\phi}}
=
\frac{
n_t(r)\,r^2\dot{\phi}
}{
\sqrt{B(r)\dot{r}^2+r^2\dot{\phi}^2}
}
=
J_t.
\label{eq:Jt_def}
\end{equation}

Let $\alpha$ be the angle between the curve and the radial direction,
measured with respect to the spatial metric $d\ell^2$. Then
$\sin\alpha = r\,d\phi/d\ell$,
and Eq.~\eqref{eq:Jt_def} becomes the generalized Snell law
\begin{equation}
n_t(r)\,r\sin\alpha
=
J_t.
\label{eq:snell_t}
\end{equation}
At the minimum radius $r=r_*$, the curve is tangent to a circle of
constant $r$, so $\alpha(r_*)=\pi/2$ and
\begin{equation}
J_t
=
n_t(r_*)\,r_*.
\label{eq:Jt_turning}
\end{equation}
Combining the Snell law with the spatial line element gives the
differential equation for the coordinate-time brachistochrone:
\begin{equation}
\frac{d\phi}{dr}
=
\frac{
n_t(r_*)\,r_*\sqrt{B(r)}
}{
r\sqrt{
n_t^2(r)\,r^2
-
n_t^2(r_*)\,r_*^2
}
}.
\label{eq:dphidr_t}
\end{equation}
For two surface points separated by an angle $\Delta$, symmetry gives
\begin{equation}
\frac{\Delta}{2}
=
\int_{r_*}^{R}
\frac{
n_t(r_*)\,r_*\sqrt{B(r)}
}{
r\sqrt{
n_t^2(r)\,r^2
-
n_t^2(r_*)\,r_*^2
}
}
\,dr.
\label{eq:opening_t}
\end{equation}
This equation determines the turning radius $r_*$ for a prescribed
opening angle $\Delta$.
Figure~\ref{fig:turning} shows the dimensionless turning radii $q^{(t)}\equiv r_{*,t}/R$ and $q^{(\tau)}\equiv r_{*,\tau}/R$ as functions of the opening angle $\Delta$, obtained by solving Eqs.~\eqref{eq:opening_t} and~\eqref{eq:opening_tau} numerically for a $\Gamma=2$ polytrope with $\mu\approx0.18$.
\begin{figure}[htb]
    \centering
    \includegraphics[width=\columnwidth]{Immagini/fig_04_turning_radius_vs_delta.png}
    \caption{Dimensionless turning radii $r_*/R$ as functions of the
    opening angle $\Delta$ for a $\Gamma=2$ polytropic star with
    $\mu\approx0.18$. Solid curve: coordinate-time brachistochrone
    $q^{(t)}$; dashed curve: proper-time brachistochrone $q^{(\tau)}$.
    Both radii decrease monotonically from $1$ (surface) at
    $\Delta\to0$ to $0$ (centre) at $\Delta\to180^\circ$. The
    coordinate-time curve lies above the proper-time curve at all
    angles, consistent with the stronger gravitational-redshift
    weighting of $n_t$ discussed in Sec.~\ref{sec:indices}.}
    \label{fig:turning}
\end{figure}
The total coordinate time is
\begin{equation}
t(\Delta)
=
\frac{2\sqrt{A_0}}{c}
\int_{r_*}^{R}
\frac{
n_t^2(r)\,r\sqrt{B(r)}
}{
\sqrt{
n_t^2(r)\,r^2
-
n_t^2(r_*)\,r_*^2
}
}
\,dr.
\label{eq:time_t}
\end{equation}

\section{Proper-time brachistochrone}

We now minimize the proper time measured along the particle worldline.
From Eq.~\eqref{eq:dtau_dl} and using
\begin{equation}
\gamma(r)=\sqrt{\frac{A_0}{A(r)}},
\qquad
v(r)=c\sqrt{1-\frac{A(r)}{A_0}},
\end{equation}
we obtain
\begin{equation}
d\tau
=
\frac{1}{c}
\sqrt{
\frac{A(r)}{A_0-A(r)}
}
\,d\ell.
\end{equation}
Thus the proper-time functional is
\begin{equation}
\tau[\gamma]
=
\frac{1}{c}
\int_\gamma
n_\tau(r)\,d\ell,
\label{eq:tau_functional}
\end{equation}
where
\begin{equation}
n_\tau(r)
=
\sqrt{
\frac{A(r)}{A_0-A(r)}
}.
\label{eq:ntau_def}
\end{equation}
The Lagrangian is
\begin{equation}
\mathcal{L}_\tau
=
n_\tau(r)
\sqrt{
B(r)\dot{r}^2+r^2\dot{\phi}^2
}.
\end{equation}
Again $\phi$ is cyclic, so
\begin{equation}
n_\tau(r)\,r\sin\alpha
=
J_\tau,
\label{eq:snell_tau}
\end{equation}
with $J_\tau = n_\tau(r_*)\,r_*$ at the turning point. Therefore
the proper-time brachistochrone satisfies
\begin{equation}
\frac{d\phi}{dr}
=
\frac{
n_\tau(r_*)\,r_*\sqrt{B(r)}
}{
r\sqrt{
n_\tau^2(r)\,r^2
-
n_\tau^2(r_*)\,r_*^2
}
}.
\label{eq:dphidr_tau}
\end{equation}
The opening-angle condition is
\begin{equation}
\frac{\Delta}{2}
=
\int_{r_*}^{R}
\frac{
n_\tau(r_*)\,r_*\sqrt{B(r)}
}{
r\sqrt{
n_\tau^2(r)\,r^2
-
n_\tau^2(r_*)\,r_*^2
}
}
\,dr,
\label{eq:opening_tau}
\end{equation}
and the total proper time is
\begin{equation}
\tau(\Delta)
=
\frac{2}{c}
\int_{r_*}^{R}
\frac{
n_\tau^2(r)\,r\sqrt{B(r)}
}{
\sqrt{
n_\tau^2(r)\,r^2
-
n_\tau^2(r_*)\,r_*^2
}
}
\,dr.
\label{eq:time_tau}
\end{equation}

Equations~\eqref{eq:opening_t}--\eqref{eq:time_t} and
\eqref{eq:opening_tau}--\eqref{eq:time_tau} define two distinct
brachistochrone problems. In general the minimizing curve for $t$ is
not the same as the minimizing curve for $\tau$, since the effective
indices $n_t$ and $n_\tau$ are different.

\section{Comparison between the two effective indices}
\label{sec:indices}

The coordinate-time and proper-time indices are
\begin{equation}
n_t(r)
=
\frac{1}
{\sqrt{A(r)\left[A_0-A(r)\right]}},
\end{equation}
and
\begin{equation}
n_\tau(r)
=
\sqrt{
\frac{A(r)}{A_0-A(r)}
}.
\end{equation}
Their ratio is
\begin{equation}
\frac{n_t(r)}{n_\tau(r)}
=
\frac{1}{A(r)}.
\end{equation}
Thus the coordinate-time problem weights the gravitational redshift
more strongly than the proper-time problem. Regions of smaller lapse
$A(r)$ are therefore favored more strongly in the minimization of $t$
than in the minimization of $\tau$.

This distinction is physically natural. The coordinate time $t$ is
the time measured by an observer at infinity, and is sensitive to
gravitational time dilation. The proper time $\tau$ is the time
experienced by the particle, and therefore the lapse enters in a
different way.

Figure~\ref{fig:effective_indices} shows the two effective indices
$n_t(r)$ and $n_\tau(r)$ computed numerically for a representative
polytropic star, together with their ratio $n_t/n_\tau = 1/A(r)$.

\begin{figure}[h]
\centering
\includegraphics[width=\columnwidth]{Immagini/fig_02_effective_indices.png}
\caption{Effective optical indices inside a polytropic TOV star
($\Gamma=2$, $\mu\approx0.18$). Panel~(a): $\log_{10} n_t$ (solid)
and $\log_{10} n_\tau$ (dashed) as functions of $r/R$. Both indices
diverge at the surface ($r\to R$) where $A\to A_0$ and
$A_0-A\to 0$.
Panel~(b): the numerically computed ratio $n_t/n_\tau$ (solid) and the analytic $1/A(r)$ (dashed), confirming the identity.
The ratio exceeds unity throughout the interior and increases toward the centre, reflecting the stronger gravitational redshift weighting of the coordinate-time functional.}
\label{fig:effective_indices}
\end{figure}

\section{Equation of state models}

The formalism above applies to any static TOV
star~\cite{ShapiroTeukolsky,Weinberg,MisnerThorneWheeler}.
In numerical applications one may begin with a polytropic equation of state,
$p=K\rho^\Gamma$,
where $K$ is the polytropic constant and $\Gamma$ is the adiabatic
index. The rest-mass density is
$\rho=(p/K)^{1/\Gamma}$,
and a simple relativistic expression for the energy density is
\begin{equation}
\epsilon
=
\rho c^2+\frac{p}{\Gamma-1}.
\end{equation}
This gives an analytic toy model for the EOS relation $\epsilon(p)$.

More realistic neutron-star calculations may instead use tabulated equations of state~\cite{Compose,OertelReview}.
In that case one constructs an interpolation
$\epsilon=\epsilon(p)$ and inserts it into the TOV equations. The
resulting functions $p(r)$, $\epsilon(r)$, $m(r)$, $A(r)$ and $B(r)$
are then used directly in the brachistochrone integrals.

\section{Numerical strategy}

A numerical implementation proceeds as follows~\cite{SilbarReddy2004}. Choose an EOS and a
central pressure $p_c$, and integrate the TOV system outward from
$r=0$ with $m(0)=0$, $p(0)=p_c$. At each step the EOS gives
$\epsilon(r)=\epsilon[p(r)]$. The stellar radius $R$ is defined by
the first zero of the pressure, and $M=m(R)$. Impose the surface
matching condition~\eqref{eq:surface_lapse} and integrate
Eq.~\eqref{eq:tov_lapse} inward to obtain $A(r)$ throughout the
interior.

Once $A(r)$ and $B(r)$ are known, construct the two effective indices
$n_t(r)$ and $n_\tau(r)$. For a prescribed opening angle $\Delta$,
solve the turning-point equation
\begin{align}
\frac{\Delta}{2}
&=
\int_{r_*}^{R}
\frac{
n_i(r_*)\,r_*\sqrt{B(r)}
}{
r\sqrt{
n_i^2(r)\,r^2
-
n_i^2(r_*)\,r_*^2
}
}
\,dr,
\nonumber\\
&\hspace{10em}
i=t,\tau.
\label{eq:turning_general}
\end{align}
yielding two generally distinct turning radii $r_{*,t}$ and
$r_{*,\tau}$. Then compute $t(\Delta)$ and $\tau(\Delta)$ from
Eqs.~\eqref{eq:time_t} and~\eqref{eq:time_tau} with the respective
turning radii. The spatial curves are reconstructed from
\begin{align}
\phi_i(r)
&=
\int_{r_{*,i}}^{r}
\frac{
n_i(r_{*,i})\,r_{*,i}\sqrt{B(\bar{r})}
}{
\bar{r}\sqrt{
n_i^2(\bar{r})\,\bar{r}^2
-
n_i^2(r_{*,i})\,r_{*,i}^2
}
}
\,d\bar{r},
\nonumber\\
&\hspace{12em}
i=t,\tau.
\label{eq:curve_reconstruction}
\end{align}

Figure~\ref{fig:brach_curves} shows the coordinate-time and
proper-time brachistochrones computed numerically for a representative polytropic star at $\Delta = 135^\circ$.

\begin{figure}[h]
\centering
\includegraphics[width=\columnwidth]{Immagini/fig_03_brachistochrone_delta135.png}
\caption{Coordinate-time brachistochrone $\phi_t(r)$ (solid) and proper-time brachistochrone $\phi_\tau(r)$ (dashed) for a test particle released from rest at the stellar surface, with opening angle $\Delta = 135^\circ$. The star is a $\Gamma = 2$ relativistic polytrope with compactness $\mu \approx 0.18$. The grey circle marks the stellar surface $r = R$. The coordinate-time curve has a larger turning radius $r_{*,t} > r_{*,\tau}$, consistent with the stronger gravitational-redshift weighting of $n_t$ relative to $n_\tau$ discussed in Sec.~\ref{sec:indices}.}
\label{fig:brach_curves}
\end{figure}

\section{Brachistochrones in anisotropic stellar interiors}
\label{sec:anisotropic}

The framework of Secs.~IV--V applies without modification to any
static, spherically symmetric interior, including stars with
anisotropic pressure. As long as spherical symmetry is preserved, the
spacetime metric retains the form~\eqref{eq:metric_AB} and the
brachistochrone equations are unchanged; only the functions $A(r)$ and
$B(r)$ that enter the effective indices differ from the isotropic case.
Pressure anisotropy in compact stars has a long
history~\cite{Ruderman1972,HerreraSantos1997,DevGleiser2002,DevGleiser2003,RaposoEtAl2019,KumarBharti2022,Becerra2024}.
In this section we specialise to the Bowers--Liang (BL)
model~\cite{BowersLiang1974}, a simple and widely used
prescription for anisotropic stellar interiors, and compare the
resulting brachistochrones with those of the corresponding isotropic
stars.

\subsection{Bowers--Liang model}

In the Bowers--Liang model the radial and tangential pressures are
related by~\cite{BowersLiang1974,CadoganPoisson2024}
\begin{equation}
p_\perp - p_\parallel
=
\lambda_{\mathrm{BL}}\,\frac{G}{c^4}\,
\frac{(\epsilon+p_\parallel)(\epsilon+3p_\parallel)}{1-2Gm/rc^2}
\,r^2,
\label{eq:BL_ansatz}
\end{equation}
where $\epsilon$ is the local energy density, $p_\parallel\equiv p$ is the
radial pressure, $p_\perp$ is the tangential pressure, and
$\lambda_{\mathrm{BL}}$ is a dimensionless anisotropy parameter. The
sign convention is such that $\lambda_{\mathrm{BL}}>0$ corresponds to
$p_\perp>p_\parallel$ (tangential pressure dominates), while
$\lambda_{\mathrm{BL}}<0$ gives $p_\perp<p_\parallel$ (radial
pressure dominates). The isotropic limit is $\lambda_{\mathrm{BL}}=0$.

The modified hydrostatic equilibrium equation reads
\begin{multline}
\frac{dp_\parallel}{dr}
=
-\frac{G(\epsilon+p_\parallel)}{c^2r^2}
\left[m(r)+\frac{4\pi r^3p_\parallel}{c^2}\right]
\\
\times\left(1-\frac{2Gm}{rc^2}\right)^{-1}
+\frac{2}{r}(p_\perp-p_\parallel).
\label{eq:BL_TOV}
\end{multline}
The mass equation and the lapse equation are unchanged:
\begin{align}
\frac{dm}{dr}
&=
\frac{4\pi r^2}{c^2}\,\epsilon(r),
\label{eq:BL_mass}
\\
\frac{d\Phi}{dr}
&=
\frac{G\left[m(r)+4\pi r^3p_\parallel/c^2\right]}
{r^2c^2\left(1-2Gm/rc^2\right)}.
\label{eq:BL_lapse}
\end{align}
With the ansatz~\eqref{eq:BL_ansatz}, Eq.~\eqref{eq:BL_TOV} becomes a
closed ODE for $p_\parallel(r)$ once an isotropic equation of state
$\epsilon=\epsilon(p_\parallel)$ is specified.
At the centre $r\to0$ the factor $r^2/(1-2Gm/rc^2)\to0$, so $p_\perp\to p_\parallel$ and the singularity of Eq.~\eqref{eq:BL_TOV} is automatically removed.
The boundary condition $p_\parallel(R)=0$ determines the stellar surface $R$, and the exterior metric is matched to Schwarzschild as in~\eqref{eq:surface_lapse}.

\subsection{Anisotropic optical indices}

Once $A(r)=e^{2\Phi(r)}$ and $B(r)=(1-2Gm/rc^2)^{-1}$ are obtained
from the integration of Eqs.~\eqref{eq:BL_mass}--\eqref{eq:BL_lapse},
the effective optical indices are given by the same expressions as in
the isotropic case,
\begin{equation}
\begin{aligned}
n_t^{(\mathrm{BL})}(r)
&=
\frac{1}{\sqrt{A(r)\left[A_0-A(r)\right]}},
\\[4pt]
n_\tau^{(\mathrm{BL})}(r)
&=
\sqrt{\frac{A(r)}{A_0-A(r)}},
\end{aligned}
\end{equation}
with $A_0=A(R)=1-2GM/(Rc^2)$. The anisotropy parameter
$\lambda_{\mathrm{BL}}$ enters only through the modification of
$A(r)$. For $\lambda_{\mathrm{BL}}>0$ (tangential pressure dominant)
the additional tangential support keeps the lapse closer to its
surface value $A_0$, so that $A_0-A(r)$ is smaller and the effective
indices are larger at depth; the optical cost of deep penetration
increases and the brachistochrone remains closer to the surface. For
$\lambda_{\mathrm{BL}}<0$ the lapse is more strongly depressed toward
the centre, $A_0-A(r)$ is larger, the indices are smaller and the
brachistochrone descends further. These trends are quantified in
Sec.~\ref{sec:anisotropic}\,C below.

\subsection{Numerical comparison at fixed compactness}

We compare isotropic ($\lambda_{\mathrm{BL}}=0$) and anisotropic
($\lambda_{\mathrm{BL}}\neq0$) models at fixed compactness
$\mu=2GM/(Rc^2)$ and fixed polytropic equation of state
$p=K\rho^\Gamma$, with $\Gamma=2$ ($n=1$ polytrope). The central
pressure $p_c$ is adjusted at each value of $\lambda_{\mathrm{BL}}$
to maintain the same $\mu$, so that the brachistochrone comparison is
performed at fixed gravitational geometry of the exterior and the same
$A_0$.

Figure~\ref{fig:brach_aniso_curves} shows the coordinate-time
brachistochrone $\phi_t(r)$ and proper-time brachistochrone
$\phi_\tau(r)$ for $\Delta=\pi/2$ and several values of
$\lambda_{\mathrm{BL}}\in[-\lambda_{\max},+\lambda_{\max}]$ at fixed
$\mu=0.2$. The two panels illustrate how pressure anisotropy shifts
the optimal tunnel depth relative to the isotropic baseline.

\begin{figure}[ht]
\centering
\includegraphics[width=0.9\columnwidth]{Immagini/fig_brach_aniso_curves.png}
\caption{Coordinate-time (solid) and proper-time (dashed)
brachistochrones for $\Delta=\pi/2$ and $\mu=0.2$, for the
Bowers--Liang model with $\lambda_{\mathrm{BL}}\in
\{-0.4,-0.2,0,+0.2,+0.4\}$ (colour-coded from blue to red). The
isotropic case ($\lambda_{\mathrm{BL}}=0$, black) serves as the
reference. Curves are plotted in the $(x/R, y/R)$ plane. The grey
circle marks the stellar surface.}
\label{fig:brach_aniso_curves}
\end{figure}

Figure~\ref{fig:qstar_vs_lambda} shows the dimensionless turning
radius $q^{(i)}\equiv r_{*,i}/R$ for $i=t,\tau$ as a function of
$\lambda_{\mathrm{BL}}$ at fixed $\Delta=\pi/2$ and $\mu=0.2$.
A positive $\lambda_{\mathrm{BL}}$ raises both brachistochrones toward the surface (larger $q$), while a negative $\lambda_{\mathrm{BL}}$ drives them deeper (smaller $q$).
The $t$--$\tau$ separation $q^{(t)}-q^{(\tau)}$, predicted analytically by the general $O(\mu)$ formula~\eqref{eq:dq1_final} derived below, is seen to depend weakly on $\lambda_{\mathrm{BL}}$ at small anisotropy but may grow or diminish at large $|\lambda_{\mathrm{BL}}|$ depending on the sign.

\begin{figure}[ht]
\centering
\includegraphics[width=0.9\columnwidth]{Immagini/fig_qstar_vs_lambda.png}
\caption{Dimensionless turning radii $q^{(t)}$ (solid) and
$q^{(\tau)}$ (dashed) as functions of $\lambda_{\mathrm{BL}}$ for
$\Delta=\pi/2$ and compactness $\mu=0.2$. The shaded band shows
the $t$--$\tau$ splitting $q^{(t)}-q^{(\tau)}$. The vertical
dashed line marks the isotropic limit $\lambda_{\mathrm{BL}}=0$.}
\label{fig:qstar_vs_lambda}
\end{figure}

The qualitative behaviour can be understood from the optical-index perspective.
For $\lambda_{\mathrm{BL}}>0$, the extra tangential support partially offsets gravitational compression, keeping the lapse $A(r)$ closer to its surface value $A_0$ throughout the interior.
The difference $A_0-A(r)$ is therefore smaller, making the effective index $n_t = [A(A_0-A)]^{-1/2}$ larger at depth: the optical cost of deep penetration increases, and the brachistochrone remains closer to the surface (larger $q$). For $\lambda_{\mathrm{BL}}<0$, the lapse is more strongly depressed toward the centre, $A_0-A(r)$ is larger, the optical index is smaller at depth, and the brachistochrone descends further (smaller $q$).
The same reasoning applies to $n_\tau$, so both turning radii shift in the same direction.

\section{Universal weak-field expansion and the $t$--$\tau$ splitting}
\label{sec:weak_field}

In this section we derive a model-independent expansion of the TOV
brachistochrone problem in the global compactness parameter
$\mu\equiv r_s/R=2GM/(Rc^2)$, keeping the Newtonian density profile
$\bar\rho(x)$ (with $x=r/R$) as an arbitrary input. The result is exact
to $O(\mu)$ for any equation of state and reduces to the
constant-density (interior Schwarzschild) case when $\bar\rho(x)=
\mathrm{const}$. While the individual turning radii $q^{(t)}\equiv
r_{*,t}/R$ and $q^{(\tau)}\equiv r_{*,\tau}/R$ separately require the
$O(\mu^2)$ metric coefficient $a_2(x)$ (and hence the pressure profile),
their \emph{difference} $q^{(t)}-q^{(\tau)}$ is fully determined already
at $O(\mu)$ and is given in closed form, modulo a single one-dimensional
quadrature, by Eq.~\eqref{eq:dq1_final} below.

\subsection{Weak-field metric coefficients}

Let $\bar m(x)$ denote the dimensionless enclosed-mass function,
\begin{equation}
\bar m(x)
=
\frac{\displaystyle\int_0^x \bar\rho(x')\,x'^2\,dx'}
     {\displaystyle\int_0^1 \bar\rho(x')\,x'^2\,dx'},
\qquad
\bar m(0)=0,\quad \bar m(1)=1.
\label{eq:mbar_def}
\end{equation}
Expanding the metric functions as
\begin{equation}
A(r)=1+\mu\,a_1(x)+O(\mu^2),
\qquad
B(r)=1+\mu\,b_1(x)+O(\mu^2),
\label{eq:AB_expansion}
\end{equation}
the exact relation $B=(1-\mu\bar m(x)/x)^{-1}$ gives immediately
\begin{equation}
b_1(x)=\frac{\bar m(x)}{x}.
\label{eq:b1_result}
\end{equation}
The lapse equation, expanded to $O(\mu^0)$ in the pressure (the
Newtonian limit of Eq.~\eqref{eq:tov_lapse}), integrates to
\begin{equation}
a_1(x)=-\Psi(x),
\qquad
\Psi(x)\equiv 1+\int_x^1\frac{\bar m(x')}{x'^2}\,dx',
\label{eq:a1_result}
\end{equation}
where the integration constant has been fixed by the surface matching
condition $A(R)=A_0=1-\mu+O(\mu^2)$, i.e.\ $a_1(1)=-1$. Equivalently,
defining the surface-relative Newtonian potential
\begin{equation}
W(x)\equiv 1-\Psi(x)=-\int_x^1\frac{\bar m(x')}{x'^2}\,dx'\;\geq 0,
\qquad W(1)=0,
\label{eq:W_def}
\end{equation}
one has $A_0-A(r)=\mu\,W(x)+O(\mu^2)$.

For a constant-density star, $\bar m(x)=x^3$, and
Eqs.~\eqref{eq:a1_result}--\eqref{eq:W_def} give
$\Psi(x)=(3-x^2)/2$ and $W(x)=(1-x^2)/2$, reproducing the leading-order
expansion of the exact interior Schwarzschild lapse.

\subsection{Leading-order optical index and Newtonian limit}

Substituting Eqs.~\eqref{eq:AB_expansion}--\eqref{eq:W_def} into the
definitions~\eqref{eq:nt_def} and~\eqref{eq:ntau_def},
\begin{multline}
n_t(x)=n_\tau(x)
=
\frac{1}{\sqrt{\mu}}\,\frac{1}{\sqrt{W(x)}}
\left[1+O(\mu)\right]
\\
\equiv
\frac{1}{\sqrt{\mu}}\,n^{(0)}(x)\left[1+O(\mu)\right],
\label{eq:n0_def}
\end{multline}
so that at leading order the two effective indices coincide, and the
overall factor $\mu^{-1/2}$ does not affect the minimizing curve. It is
convenient to introduce the leading-order optical impact function
\begin{equation}
P_0(x)\equiv n^{(0)}(x)\,x=\frac{x}{\sqrt{W(x)}},
\qquad
P_0^2(x)=\frac{x^2}{W(x)}.
\label{eq:P0_def}
\end{equation}
The Newtonian Snell law and opening-angle condition then read
\begin{equation}
P_0(x)\sin\alpha=P_0(q_0),
\label{eq:snell0}
\end{equation}
\begin{equation}
\frac{\Delta}{2}
=
\int_{q_0}^1
\mathcal{I}(x;q_0)\,dx
\equiv
\mathcal{F}^{(0)}[q_0],
\label{eq:F0_def}
\end{equation}
where the Newtonian kernel is
\begin{equation}
\mathcal{I}(x;q)
\equiv
\frac{P_0(q)}
{x\sqrt{P_0^2(x)-P_0^2(q)}}
=
\frac{q/\sqrt{W(q)}}
{x\sqrt{\dfrac{x^2}{W(x)}-\dfrac{q^2}{W(q)}}}.
\label{eq:kernel_def}
\end{equation}
Equation~\eqref{eq:F0_def} determines $q_0(\Delta)$ for an arbitrary
density profile $\bar\rho(x)$ via $W(x)$, generalizing
$q_0=1-\Delta/\pi$ (the value recovered for $W(x)=(1-x^2)/2$, for which
$\mathcal{F}^{(0)}[q_0]=\tfrac{\pi}{2}(1-q_0)$).

\subsection{First-order splitting of the optical indices}

At the next order, write
\begin{equation}
n_i(x)
=
\frac{1}{\sqrt{\mu}}\,n^{(0)}(x)
\left[1+\mu\,\delta_i(x)+O(\mu^2)\right],
\qquad i=t,\tau.
\label{eq:ni_expansion}
\end{equation}
The exact ratio $n_t/n_\tau=1/A(r)$ gives, using
$A^{-1}(r)=1+\mu\Psi(x)+O(\mu^2)$,
\begin{equation}
\delta_t(x)-\delta_\tau(x)=\Psi(x).
\label{eq:delta_diff}
\end{equation}
Equation~\eqref{eq:delta_diff} is exact at this order and depends only
on $\Psi(x)$ from Eq.~\eqref{eq:a1_result}; in particular it does not
require the $O(\mu^2)$ lapse coefficient $a_2(x)$, which would instead
be needed for $\delta_\tau(x)$ individually.

\subsection{Perturbative inversion of the opening-angle condition}

Writing $q^{(i)}=q_0+\mu\,q_1^{(i)}+O(\mu^2)$ and expanding the
opening-angle integral~\eqref{eq:opening_t} or~\eqref{eq:opening_tau}
(generic form~\eqref{eq:turning_general}) to $O(\mu)$ at fixed $\Delta$,
the integrand receives three distinct first-order contributions: from
the index correction $\delta_i$, from the factor $\sqrt{B(x)}=1+
\tfrac{\mu}{2}b_1(x)+O(\mu^2)$, and from the resulting shift of the
turning point. The $\sqrt{B}$ contribution is common to both problems
and therefore cancels in the difference $q^{(t)}-q^{(\tau)}$; only the
$\delta_i$-dependent part survives. The inversion gives
\begin{equation}
q_1^{(i)}
=
-\,
\frac{\mathcal{G}_i[q_0]+\mathcal{B}[q_0]}
{\partial_q\mathcal{F}^{(0)}(q_0)},
\label{eq:q1_general}
\end{equation}
where $\mathcal{B}[q_0]$ collects the ($i$-independent) $\sqrt{B}$
contribution and $\mathcal{G}_i[q_0]$ is the linear functional of
$\delta_i$ obtained from the $O(\mu)$ variation of the kernel
$\mathcal{I}(x;q)$ in Eq.~\eqref{eq:kernel_def}, evaluated along the
$O(\mu^0)$ trajectory.

\subsection{The $t$--$\tau$ splitting at $O(\mu)$}

Subtracting, the common term $\mathcal{B}$ drops out and, by
linearity together with Eq.~\eqref{eq:delta_diff},
$\mathcal{G}_t[q_0]-\mathcal{G}_\tau[q_0]=\mathcal{G}[q_0]$, where
$\mathcal{G}[q_0]$ is the same functional evaluated on $\Psi(x)$ alone.
Carrying out the $O(\mu)$ variation of Eq.~\eqref{eq:kernel_def}
explicitly and simplifying, the bracket collapses to a single term and
\begin{equation}
\mathcal{G}[q_0]
=
\int_{q_0}^1
\mathcal{I}(x;q_0)\,
\frac{P_0^2(x)\,\bigl[\Psi(q_0)-\Psi(x)\bigr]}
     {P_0^2(x)-P_0^2(q_0)}
\,dx,
\label{eq:G_functional}
\end{equation}
with $P_0^2(x)=x^2/W(x)$ from Eq.~\eqref{eq:P0_def} and, equivalently,
$\Psi(q_0)-\Psi(x)=W(x)-W(q_0)$. The integrand is integrable at the
turning point $x\to q_0$, where the ratio remains finite while
$\mathcal{I}\sim(x-q_0)^{-1/2}$. Combining
Eqs.~\eqref{eq:q1_general} and~\eqref{eq:G_functional}, the $t$--$\tau$
splitting of the turning radius is
\begin{equation}
\boxed{
\begin{aligned}
q^{(t)}-q^{(\tau)}
&=
\mu\left(q_1^{(t)}-q_1^{(\tau)}\right)+O(\mu^2)
\\[4pt]
&=
-\,\mu\,
\frac{\mathcal{G}[q_0]}{\partial_q\mathcal{F}^{(0)}(q_0)}
+O(\mu^2).
\end{aligned}
}
\label{eq:dq1_final}
\end{equation}

Equation~\eqref{eq:dq1_final} is the central result of this section:
for \emph{any} equation of state, encoded through the single function
$W(x)$ via Eqs.~\eqref{eq:mbar_def} and~\eqref{eq:W_def}, the leading
relativistic separation between the coordinate-time and proper-time
brachistochrones is obtained from two one-dimensional quadratures ---
$\mathcal{F}^{(0)}(q_0)$ and $\mathcal{G}[q_0]$ --- and one derivative,
without solving the full TOV system.

Figure~\ref{fig:splitting} shows $q^{(t)}-q^{(\tau)}$ as a function of $\mu$ for polytropic equations of state with $\Gamma=2$, $3$, and $10$, computed numerically from the full TOV solutions (solid curves), together with the $O(\mu)$ weak-field prediction of Eq.~\eqref{eq:dq1_final} (dashed curves).
The analytical Schwarzschild result~\eqref{eq:dq_schw} is shown for reference (dash-dotted).
Over the compactness range shown, the weak-field prediction reproduces the trend of the full numerical result, with the deviation growing as $\mu$ increases beyond $\mu\gtrsim0.15$, as expected for an $O(\mu)$ expansion.

\begin{figure}[htb]
    \centering
    \includegraphics[width=\columnwidth]{Immagini/fig_14_t_tau_splitting_o1.png}
    \caption{Difference $q^{(t)}-q^{(\tau)}$ between the coordinate-time
    and proper-time brachistochrone turning radii as a function of the
    compactness $\mu=r_s/R$, for polytropic equations of state with
    $\Gamma=2$, $3$, $10$ (solid curves). Dashed curves show the
    weak-field prediction to $O(\mu)$ from Eq.~\eqref{eq:dq1_final};
    the dash-dotted line is the Schwarzschild analytical
    result~\eqref{eq:dq_schw}. All curves are for $\Delta=135^\circ$.}
    \label{fig:splitting}
\end{figure}

\subsection{Constant-density limit and closed form}

For $\bar\rho(x)=\mathrm{const}$, $W(x)=(1-x^2)/2$ and
$\Psi(x)=(3-x^2)/2$, so that $P_0^2(x)=2x^2/(1-x^2)$ and the kernel
reduces to $\mathcal{I}(x;q)=q\sqrt{1-x^2}\,/\,(x\sqrt{x^2-q^2})$. The
integrand of Eq.~\eqref{eq:G_functional} then collapses to the
elementary form
\begin{equation}
\mathcal{I}(x;q)\,
\frac{P_0^2(x)\bigl[\Psi(q)-\Psi(x)\bigr]}{P_0^2(x)-P_0^2(q)}
=
\frac{q(1-q^2)}{2}\,
\frac{x\sqrt{1-x^2}}{\sqrt{x^2-q^2}},
\label{eq:integrand_schw}
\end{equation}
and, using $\int_{q}^{1}x\sqrt{1-x^2}/\sqrt{x^2-q^2}\,dx=
\pi(1-q^2)/4$, one finds the closed form
\begin{equation}
\mathcal{G}[q_0]=\frac{\pi\,q_0\,(1-q_0^2)^2}{8}.
\label{eq:G_schw}
\end{equation}
Since $\mathcal{F}^{(0)}[q_0]=\tfrac{\pi}{2}(1-q_0)$ gives
$\partial_q\mathcal{F}^{(0)}=-\pi/2$, Eq.~\eqref{eq:dq1_final} yields
\begin{equation}
q^{(t)}-q^{(\tau)}
=
\mu\,\frac{q_0\,(1-q_0^2)^2}{4}
+O(\mu^2).
\label{eq:dq_schw}
\end{equation}
The coefficient is strictly positive for $0<q_0<1$, confirming that the
coordinate-time brachistochrone is shallower (larger turning radius)
than the proper-time one, consistently with the discussion of
Sec.~\ref{sec:weak_field}. For $\Delta=\pi/2$ ($q_0=1/2$) the
coefficient equals $9/128$. This analytic result provides a direct
check of the general formula~\eqref{eq:dq1_final} and a closed-form
benchmark against which the numerical TOV integration of arbitrary
equations of state can be validated. It is also consistent with the
NLO coordinate-time coefficient
$q_1^{(t)}=q_0(1-q_0^2)(15-7q_0^2)/32$ obtained for the interior
Schwarzschild metric in Ref.~\cite{Rosignoli2026}, of which
Eq.~\eqref{eq:dq_schw} represents the proper-time counterpart implicit
in that analysis.

Figure~\ref{fig:schw_validation} shows $q^{(t)}$ (top) and $q^{(\tau)}$ (bottom) as functions of the compactness $\mu$ for polytropic equations of state with $\Gamma=5$, $10$, $20$, together with the incompressible ($\rho=\mathrm{const}$) Schwarzschild reference and the Newtonian limit $q_0=1-\Delta/\pi$.
As $\Gamma$ increases the stellar density profile becomes more centrally peaked and the polytropic curves approach the Schwarzschild reference from above for $q^{(t)}$ and from below for $q^{(\tau)}$, consistently with the stronger gravitational redshift weighting of $n_t$ relative to $n_\tau$.

\begin{figure}[htb]
    \centering
    \includegraphics[width=\columnwidth]{Immagini/fig_15_schwarzschild_validation.png}
    \caption{Dimensionless turning radii $q^{(t)}$ (top) and
    $q^{(\tau)}$ (bottom) as functions of the compactness
    $\mu = r_s/R$, for polytropic equations of state with
    $\Gamma = 5$, $10$, $20$ (coloured curves), the incompressible
    Schwarzschild interior ($\rho = \mathrm{const}$, black), and the
    Newtonian limit (grey dashed). All curves are for
    $\Delta = 135^\circ$. The polytropic curves converge toward the
    Schwarzschild reference as $\Gamma\to\infty$, validating the
    analytical result~\eqref{eq:dq_schw}.}
    \label{fig:schw_validation}
\end{figure}

\section{Expected observables}

The main observables of the problem are: (1)~the brachistochrone
curves $\phi_t(r)$ and $\phi_\tau(r)$; (2)~the turning radii
$r_{*,t}$ and $r_{*,\tau}$; (3)~the travel times $t(\Delta)$ and
$\tau(\Delta)$; (4)~their dependence on the compactness
$C=GM/(Rc^2)$; and (5)~their dependence on the equation of state.

A useful comparison is obtained by studying the sequence
\begin{align}
&\text{Newtonian}
\;\longrightarrow\;
\text{int.\ Schwarzschild}
\nonumber\\
&\quad\longrightarrow\;
\text{TOV polytrope}
\;\longrightarrow\;
\text{realistic EOS},
\label{eq:sequence}
\end{align}
which isolates the successive effects of relativistic redshift,
spatial curvature and realistic density stratification.

\section{Discussion}

Figure~\ref{fig:eos_comparison} illustrates the last step of this sequence: the coordinate-time and proper-time brachistochrones computed for a $\Gamma=2$ relativistic polytrope and for a realistic tabulated equation of state (CompOSE \texttt{qmc\_rmf\_1}), both matched to the same compactness $\mu\approx0.18$ and opening angle $\Delta=135^\circ$.
The spatial curves are nearly indistinguishable at the scale of the stellar radius, but the inset reveals a systematic shift in the turning radii: the realistic EOS, having a different density stratification, places both brachistochrones at slightly different depths than the polytrope at the same $\mu$.

\begin{figure}[htb]
    \centering
    \includegraphics[width=\columnwidth]{Immagini/fig_07_eos_comparison.png}
    \caption{Coordinate-time (solid/dash-dotted) and proper-time
    (dashed/dotted) brachistochrones in the $(x/R,\,y/R)$ plane for a
    $\Gamma=2$ polytrope (blue/orange) and a realistic CompOSE EOS
    \texttt{qmc\_rmf\_1} (green/yellow), both at compactness
    $\mu\approx0.18$ and opening angle $\Delta=135^\circ$. The grey
    dotted circle marks the stellar surface. The inset zooms on the
    turning-point region, showing that the two equations of state
    predict slightly different values of $q^{(t)}$ and $q^{(\tau)}$
    at the same compactness.}
    \label{fig:eos_comparison}
\end{figure}

The TOV brachistochrone problem provides a geometrical probe of the
interior structure of a compact star. The effective indices $n_t$ and
$n_\tau$ are determined by the lapse function $A(r)$, while the
spatial curvature enters through $B(r)$. Both functions are fixed by
the equation of state through the TOV equations.

Since $n_t(r)=n_\tau(r)/A(r)$ and $A(r)\leq A_0<1$, the
coordinate-time index $n_t$ is everywhere larger than $n_\tau$.
For fixed surface aperture $\Delta$, the turning-point integral
saturates at a larger $r_*$ for the coordinate-time problem,
meaning the coordinate-time brachistochrone remains closer to
the stellar surface than the proper-time brachistochrone.
Equivalently, for fixed turning radius $r_*$, the
coordinate-time curve subtends a larger aperture than the
proper-time curve. This is consistent with the flat-spacetime (Rindler) analogue,
where uniform acceleration plays the role of the lapse,
and the coordinate-time brachistochrone is stretched relative
to the Newtonian cycloid while the proper-time brachistochrone
is compressed.


The present treatment assumes a static, spherically symmetric star and
a test particle moving without backreaction. It neglects rotation,
magnetic fields, viscosity, crustal effects and hydrodynamical
interaction with the stellar matter. Future extensions could include
slowly rotating Hartle--Thorne interiors, anisotropic pressure,
charged fluids or modified gravity stellar models.

\section{Conclusion}

We have derived the coordinate-time and proper-time brachistochrone
equations for a massive test particle moving inside a static TOV
stellar interior. Given an equation of state, the TOV equations
determine the metric functions $A(r)$ and $B(r)$. The brachistochrone
problem then reduces to two Fermat-type variational principles with
effective indices
\begin{equation}
n_t(r)
=
\frac{1}
{\sqrt{
A(r)\left[A_0-A(r)\right]
}},
\end{equation}
and
\begin{equation}
n_\tau(r)
=
\sqrt{
\frac{A(r)}{A_0-A(r)}
}.
\end{equation}
The angular symmetry yields generalized Snell laws,
\begin{equation}
n_i(r)\,r\sin\alpha
=
n_i(r_*)\,r_*,
\qquad
i=t,\tau,
\end{equation}
from which the turning radius, the curve and the travel time are
obtained by quadrature.

This framework generalizes the relativistic gravitational-train problem
from constant-density interiors to realistic stellar models generated
by equations of state. The main analytical result is
Eq.~\eqref{eq:dq1_final}: for \emph{any} equation of state, encoded
in the single function $W(x)$ through the Newtonian density profile,
the leading relativistic separation $q^{(t)}-q^{(\tau)}$ between the
two brachistochrones is obtained from two one-dimensional
quadratures and one derivative, bypassing the full TOV integration.
The constant-density limit yields the closed form $\mu\,q_0(1-q_0^2)^2/4$
(Eq.~\eqref{eq:dq_schw}), which provides an analytical benchmark
against which numerical solutions for arbitrary equations of state can
be validated, and which extends the NLO perturbative analysis of
Ref.~\cite{Rosignoli2026} to the proper-time brachistochrone.
This opens the possibility of using the $t$--$\tau$ splitting as a
geometrical probe of the density stratification of compact stars,
complementary to standard mass-radius observations.

\section*{Data Availability}
The numerical code used to produce the figures in this paper
is publicly available at
\url{https://github.com/ImanetorX98/TOVBrachistochrone}.
The CompOSE equation-of-state tables used in Sec.~XII are
available from the CompOSE database at
\url{https://compose.obspm.fr}~\cite{Compose}.

\begin{acknowledgments}
The author declares no conflicts of interest.
The author acknowledges the University of Pavia for the academic
training and scientific environment that supported the development of
this work. The author also acknowledges the use of AI-based tools,
namely OpenAI’s GPT-5.5 Thinking and Anthropic’s Claude Sonnet 4.6,
for assistance in preparing an initial draft of parts of the derivation
and in supporting the editing of the manuscript. The author reviewed
and verified all AI-assisted content and retains full responsibility
for the scientific content, calculations, and conclusions of this work.

No funds, grants, or other support were received during the preparation
of this manuscript.
\end{acknowledgments}

\bibliographystyle{apsrev4-2}

\begin{thebibliography}{99}

\bibitem{Klotz2015}
A.~R. Klotz,
``The Gravity Tunnel in a Non-Uniform Earth,''
Am.\ J.\ Phys.\ \textbf{83}, 231 (2015).

\bibitem{GjerlovPesnell2019}
A.~Gjerl\o{}v and W.~D. Pesnell,
``Orbits through polytropes,''
Am.\ J.\ Phys.\ \textbf{87}, 452 (2019).

\bibitem{Parker2017}
E.~Parker,
``A relativistic gravity train,''
Gen.\ Relativ.\ Gravit.\ \textbf{49}, 106 (2017).

\bibitem{Seel2018}
M.~Seel,
``The relativistic gravity train,''
Eur.\ J.\ Phys.\ \textbf{39}, 035603 (2018).

\bibitem{Perlick1990}
V.~Perlick,
``On Fermat's Principle in General Relativity. I. The General Case,''
Class.\ Quantum Grav.\ \textbf{7}, 1319 (1990).

\bibitem{Gibbons2016}
G.~W. Gibbons,
``The Jacobi-metric for timelike geodesics in static spacetimes,''
Class.\ Quantum Grav.\ \textbf{33}, 025004 (2016).

\bibitem{Schmidt2018}
H.-J. Schmidt,
``Snellius meets Schwarzschild: Refraction of brachistochrones and
time-like geodesics,''
arXiv:1809.00355 [gr-qc] (2018).

\bibitem{Rosignoli2026}
I.~Rosignoli,
``The Gravitational Train in the Interior Schwarzschild Metric:
NLO and NNLO Perturbative Expansions,''
preprint (2026), Zenodo:
\href{https://doi.org/10.5281/zenodo.20617416}{10.5281/zenodo.20617416}.

\bibitem{Tolman1939}
R.~C. Tolman,
``Static Solutions of Einstein's Field Equations for Spheres of Fluid,''
Phys.\ Rev.\ \textbf{55}, 364 (1939).

\bibitem{OppenheimerVolkoff1939}
J.~R. Oppenheimer and G.~M. Volkoff,
``On Massive Neutron Cores,''
Phys.\ Rev.\ \textbf{55}, 374 (1939).

\bibitem{ShapiroTeukolsky}
S.~L. Shapiro and S.~A. Teukolsky,
\textit{Black Holes, White Dwarfs, and Neutron Stars:
The Physics of Compact Objects}
(Wiley, New York, 1983).

\bibitem{Weinberg}
S.~Weinberg,
\textit{Gravitation and Cosmology}
(Wiley, New York, 1972).

\bibitem{MisnerThorneWheeler}
C.~W. Misner, K.~S. Thorne, and J.~A. Wheeler,
\textit{Gravitation}
(Freeman, San Francisco, 1973).

\bibitem{Compose}
S.~Typel, M.~Oertel, and T.~Kl\"{a}hn,
``{CompOSE} -- {CompStar} Online Supernovae Equations of State,''
Phys.\ Part.\ Nucl.\ \textbf{46}, 633 (2015).

\bibitem{OertelReview}
M.~Oertel, M.~Hempel, T.~Kl\"{a}hn, and S.~Typel,
``Equations of state for supernovae and compact stars,''
Rev.\ Mod.\ Phys.\ \textbf{89}, 015007 (2017).

\bibitem{SilbarReddy2004}
R.~R. Silbar and S.~Reddy,
``Neutron Stars for Undergraduates,''
Am.\ J.\ Phys.\ \textbf{72}, 892 (2004).

\bibitem{Ruderman1972}
M.~Ruderman,
``Pulsars: Structure and Dynamics,''
Ann.\ Rev.\ Astron.\ Astrophys.\ \textbf{10}, 427 (1972).

\bibitem{HerreraSantos1997}
L.~Herrera and N.~O. Santos,
``Local Anisotropy in Self-Gravitating Systems,''
Phys.\ Rep.\ \textbf{286}, 53 (1997).

\bibitem{DevGleiser2002}
K.~Dev and M.~Gleiser,
``Anisotropic Stars: Exact Solutions,''
Gen.\ Relativ.\ Gravit.\ \textbf{34}, 1793 (2002).

\bibitem{DevGleiser2003}
K.~Dev and M.~Gleiser,
``Anisotropic Stars II: Stability,''
Gen.\ Relativ.\ Gravit.\ \textbf{35}, 1435 (2003).

\bibitem{RaposoEtAl2019}
G.~Raposo, P.~Pani, M.~Bezares, C.~Palenzuela, and V.~Cardoso,
``Anisotropic Stars as Ultracompact Objects in General Relativity,''
Phys.\ Rev.\ D \textbf{99}, 104072 (2019).

\bibitem{KumarBharti2022}
J.~Kumar and P.~Bharti,
``Relativistic models for anisotropic compact stars: A review,''
New Astron.\ Rev.\ \textbf{95}, 101662 (2022).

\bibitem{Becerra2024}
L.~M. Becerra, E.~A. Becerra-Vergara, and F.~D. Lora-Clavijo,
``Realistic anisotropic neutron stars: Pressure effects,''
Phys.\ Rev.\ D \textbf{109}, 043025 (2024).

\bibitem{BowersLiang1974}
R.~L. Bowers and E.~P.~T. Liang,
``Anisotropic Spheres in General Relativity,''
Astrophys.\ J.\ \textbf{188}, 657 (1974).

\bibitem{CadoganPoisson2024}
T.~Cadogan and E.~Poisson,
``Self-gravitating anisotropic fluids. I: Context and overview,''
arXiv:2406.03185 [gr-qc] (2024).

\end{thebibliography}

\end{document}
